
% !TEX root = ../../../Esperimenti/.tex/MEF.tex
% LTeX: language=it

\nodeDiagramSettings
\pgfmathsetmacro\lengthAxis{1.5}
\pgfmathsetmacro\lengthPlot{\lengthAxis-.25}
\tikzset{/pgf/fpu/install only={reciprocal,exp,ln}}

\tikzset{
    distributionStyle/.style={
        color=black!50,
        opacity=.7,
        % color=black,
        % draw=none,
        % line width=.5pt,
    },
    plotDomain/.style={
        domain=-\lengthPlot:\lengthPlot,
        samples=300,
        smooth,
    }
}

\begin{tikzpicture}[
    declare function={
        gauss(\x,\m,\s,\a)=\a*exp(-((\x-\m)^2)/(2*\s^2)); % Funzione gaussiana
        beta(\x,\l,\a,\b,\c)=\c*((\x+\l)/(2*\l))^(\a-1)*(1-(\x+\l)/(2*\l))^(\b-1); % Funzione beta nell'intervallo [-\l,\l]
    },
]%
    %\begingroup Coordinate dei vertici
        \newcommand\vertexCoordinates{
            % Punti interni visibili
            \coordinate (i) at (-4,0);
            \coordinate (k) at (-.5,1);
            \coordinate (j) at (2,-2);

            \coordinate (ai) at ([xshift=-2cm,yshift=-1cm]i);
            \coordinate (ak) at ([xshift=-.5cm,yshift=-2cm]k);
            \coordinate (aj) at ([xshift=2cm,yshift=1cm]j);
        }

        \newcommand\drawNodes{
            \foreach \i in {i,k,j}
                \node[
                    nodeStyle,
                    numberStyle,
                    fill=white
                ] (n\i) at (\i) {\i};
        }

        \newcommand\drawEdges[6]{
            % #1 -> Primo nodo
            % #2 -> Punto medio dei lati entranti e uscenti del primo nodo
            % #1 -> Secondo nodo
            % #4 -> Punto medio dei lati entranti e uscenti del secondo nodo
            % #5 -> Distanza del punto di contatto dei lati rispetto a quello medio
            % Tutt'i punti sono posizionati lungo la circonferenza per cui i loro valori vanno da 0 a 360
            %
            \draw[nodeArrowExact] (#1) to[out=#2+#5,in=#4-#6] (#3);
            \draw[nodeArrowExact] (#3) to[out=#4+#5,in=#2-#6] (#1);
        }

        \newcommand\drawAxes{
            \foreach \i in {i,k,j}{
                \draw[arrowAxis] 
                    ([xshift=-\lengthAxis cm]a\i) -- ([xshift=\lengthAxis cm]a\i)
                    node[commentStyle,below right=0pt] {\(s\)};
            }
        }

        \newcommand\DrawDistributions{
            % % La chiave «domain» può accettare unità di misura (tipo «cm», «pt» o «em») ma le tratta in genere molto male trasformandole in una quantità molto piú grande di quanto scritto; pertanto è meglio usare numeri «senza» unità di misura che verranno interpretati rispetto all'unità di base del gruppo considerato (di solito è in centrimetri)
 
            \begin{scope}[shift={(ai)}]
                \fill[distributionStyle] plot[plotDomain]
                    (\x,{.5*beta(\x,\lengthPlot,2,8,45)+.5*beta(\x,\lengthPlot,3.5,3,4)}) -- 
                    % (\x,{gauss(\x,0,.4,1)}) -- 
                    (\lengthAxis cm,0) -- (-\lengthAxis cm,0) -- cycle;
                \node[commentStyle,above] at (.5,.75) {$f_i(s,t)$};
            \end{scope} % Distribuzione i: Gaussiana singola centrata

            \begin{scope}[shift={(ak)}]
                \fill[distributionStyle] plot[plotDomain]
                    % (\x,{.5*beta(\x,\lengthPlot,2,6,0.8)+.5*beta(\x,\lengthPlot,6,2,0.4)}) --
                    (\x,{gauss(\x,-.5*\lengthPlot,.2,1.2) + gauss(\x,.25*\lengthPlot,.3,.75)}) -- 
                    (\lengthAxis cm,0) -- (-\lengthAxis cm,0) -- cycle;
                \node[commentStyle,above] at (.25,.75) {$f_k(s,t)$};
            \end{scope} % Distribuzione i: Gaussiana singola centrata

            \begin{scope}[shift={(aj)}]
                \fill[distributionStyle] plot[plotDomain]
                    (\x,{beta(\x,\lengthPlot,2,2,3)}) --
                    % (\x,{gauss(\x,.3,.35,.8)}) -- 
                    (\lengthAxis cm,0) -- (-\lengthAxis cm,0) -- cycle;
                \node[commentStyle,above] at (-1,.75) {$f_j(s,t)$};
            \end{scope} % Distribuzione i: Gaussiana singola centrata
        }
    %\endgroup

    \begin{scope}[
        local bounding box=tikZFigure,
        scale=\plotScale,
        transform shape,
    ]
        \vertexCoordinates % Coordinate
        \drawNodes{n} % Nodi

        % Lati
        \drawEdges{ni}{-65}{nj}{175}{10}{10}
        \drawEdges{nj}{90}{nk}{-45}{20}{10}
        \drawEdges{nk}{190}{ni}{45}{10}{10}

        \DrawDistributions % Distribuzioni della taglia
        \drawAxes % Assi della taglia

        %\begingroup Frecce
            % \draw[transformingArrow] 
            %     ($(ExactGraph.south west)!0.2!(ExactGraph.south east)$) 
            %     to[bend right=40] node[arrowCommentStyle] {\cref{defMatriceGradi}} (ApproxGraph.west);

            % \draw[transformingArrow]
            %     (ApproxGraph.east) to[bend right=40] node[arrowCommentStyle] {\cref{teoRilassamentoTopologia}}
            %     ($(RelaxedGraph.south west)!0.8!(RelaxedGraph.south east)$);

            % % \draw[transformingArrow] (ExactGraph.south) to[bend right=40] (ApproxGraph.west);
            % % \draw[transformingArrow] (ApproxGraph.east) to[bend right=40] (RelaxedGraph.south);

            % % \draw[transformingArrow] (ExactGraph.south) |- (ApproxGraph.west);
            % % \draw[transformingArrow] (ApproxGraph.east) -| (RelaxedGraph.south);
        %\endgroup
    \end{scope}

    % \begin{scope}[on background layer]
    %     \draw[style=help lines] (tikZFigure.north west) grid[step=1em] (tikZFigure.south east);
    %     \node at (0,0) {O};
    % \end{scope}
\end{tikzpicture}
